\newpage

\scan{122}

\newpage

% Translation below

% Går dig någon med sin halfwa styrkia innan till hos din swaga / så
% låt din udd med en hast siunka / som tillförende; Gack åter hastigt
% med din Klinga öfwer sig / och tag med din halfwa styrkia utan till
% fast i hans halfẃa swaga / string' honom hans Klinga / går han med
% en hast under dink Klinga igenom / så gif gran acht på tempo i det
% han går igenom / at du träder med din högra fot på honom in / och
% stöt med secunda innan till hans wenstra bröst.

\exercise{}
% Ligger du med en lång Klinga för din man / och du wore utan till
% string', så cavera med en hast under hans Klinga / och stöt med
% qvarta utan till öfwer hans wenstra Arm till hans bröst: Skulle han
% samma qvart' med sin Klinga till sin wenstra sijda parera, så
% passera i en hast på honom ut/ och stöt secund' utan till öfwer hans
% wenstra Arm.

\exercise{}
% Ligger någon i en jämbn tern' för dig / så string' honom hans Klinga
% utan till; Caverar han då med en hast under din Klinga / och stöter
% med qvarta utan till öfwer din högra Arm / och stöten geck
% temmeligen högt till don högre Axel / så gif noga acht i det han
% stöter / at du i det samma förfaller under hans Klinga / och stöter
% med secunda till hans wenstre bröst.

\exercise{}
% Är du string', så låt din udd siunka till jorden; Wille han då med
% sin klinga föllia din effter / går du med din Klinga under hans
% fram/ och giör honom en Fint innan till med jern / då lärer han taga
% med Klingan till sin högra sijda / i det samma caverar du under hans
% klinga och stöter med qvarta utan till öfwer hans wenstre Arm med
% sin Klinga / så passera i en hast under hans Klinga dinkos / och
% stöt med secund' utan till öfwer hans wenstre Arm.

\exercise{}
% Låt din udd åter siunka som tillförende / föllier han då med sin
% Klinga din effter / så giör honom i snarhet en Fint öfwer och under
% hans klinga; Förfarer han då med sin klinga nederåth / så gack åter
% under hans klinga igenom / och stöt med qvarta utan till ögwer hans
% wenstre Arm.

% \exercise{huru man emot Finter sig förhålla skall.}

% Stringera honom innan till; Caverar han med det samma och giör dig
% en Fint innan till utan jern / med en battute af foten / så voltera
% qvartam der emot med uddens försänkning under hans klinga / så
% blifwer honom genomgången till utwärtes stöten förtagen: Men giör
% han samma Fint med jern / så gack i det samma med din styrkia till
% hans swaga / och stöt utan till öfwer hans klinga till hans högre
% Axel; Eller parera med willia lijtet effter den förrige Finten som
% skedde utan jern / allenast med fremste leden aff wenstre handen;
% Will han på den paraden utan till öfwer din klinga instöta / så
% falzera hans klinga / och stöt utan till med qvart' under hans
% klinga / Eller giör honom en contra-fint
